\section{Experimental design}
\label{sec:design}

\subsection{Model potential and reference}
We use a two-dimensional model potential with reaction coordinate
$\xib(x,y)=x$,
\begin{equation}
  \begin{aligned}
    V(x,y) =\;& 3\,e^{-x^2}\Bigl(e^{-(y-1/3)^2}-e^{-(y-5/3)^2}\Bigr)
    \;-\;5\,e^{-y^2}\Bigl(e^{-(x-1)^2}+e^{-(x+1)^2}\Bigr)\\
    &+\;0.2\,x^4 + 0.2\,(y-1/3)^4 .
  \end{aligned}
  \label{eq:potential}
\end{equation}
As shown in \cref{fig:geometry}, $V$ has two metastable basins near $x=\pm1$
separated by a barrier near $x=0$ (with an upper sub-basin near $(0,5/3)$), so
unbiased dynamics is metastable in $x$. We fix $\beta=4$. The reference
$\Fref,\Fpref,\pref$ are obtained by the $y$-quadrature of
\eqref{eq:pref}--\eqref{eq:meanforce} on a fine grid and are treated as ground
truth.

\begin{figure}[t]
  \centering
  \includegraphics[width=0.92\textwidth]{fig01_reference_geometry.png}
  \caption{Problem geometry. (a) The potential $V(x,y)$~\eqref{eq:potential}
    with two basins near $x=\pm1$. (b) The canonical density
    $\pi\propto e^{-\beta V}$ at $\beta=4$, concentrated in the two basins.
    (c) The reference free energy $\Fref(x)$ (a double well in $x$) and
    (d) the reference mean force $\Fpref(x)$, both by $y$-quadrature.}
  \label{fig:geometry}
\end{figure}

\subsection{Methods compared}
Four methods are compared, identical except for the FR target:
\textbf{ABF only} (no birth--death; baseline);
\textbf{ABF+FR estimated}~\eqref{eq:q_est} (the practical method);
\textbf{ABF+FR uniform}~\eqref{eq:q_unif} (naive marginal-flattening ablation);
and \textbf{ABF+FR oracle}~\eqref{eq:q_oracle} (diagnostic control,
\cref{rem:oracle}). All four share the same potential, integrator, ABF
estimator, seeds and initial conditions, so differences are attributable to the
FR step and its target alone (matched-seed comparison).

\subsection{Metrics}
\label{sec:metrics}
All profile errors are interior, domain-RMS $\Ltwo$ distances on the window
$x\in[-2.5,2.5]$ (the reflecting walls and the $x^4$ confinement make the edges
uninformative), with free-energy profiles centred as in \cref{sec:abf}.

\paragraph{Primary (practical) metrics.}
\begin{equation}
  \Norm{\Fhat_T-\Fref}_{\Ltwo},\qquad
  \Norm{\Fphat_T-\Fpref}_{\Ltwo},\qquad
  \int_0^T \Norm{\Fhat_t-\Fref}_{\Ltwo}\dd t.
  \label{eq:primary}
\end{equation}
The first two are \emph{final-budget accuracy} at the end of the run; the
\emph{integrated} error rewards \emph{fast} convergence (anytime performance),
since a method that reaches a good profile earlier accumulates less error over
$[0,T]$. These two notions need not be optimised by the same hyperparameters,
a point we treat carefully in \cref{sec:results}.

\paragraph{Diagnostic metrics.}
Reaction-coordinate marginal flatness ($\Norm{\phat-\mathrm{Unif}}$ and
$\Norm{\phat-q}$); cumulative barrier crossings at $x=0$; region fractions
(left basin / barrier / right basin); the FR event fraction (particles replaced
per FR application) and the FR score dispersion $\mathrm{std}_i\,S_i$;
conditional $Y\mid X=x$ errors (\cref{sec:conditional}); and ancestor/diversity
counts. We stress that \emph{marginal flatness is a diagnostic, not the
objective}: the goal is free-energy/mean-force accuracy~\eqref{eq:primary}. A
method can flatten $\phat$ and yet fail to improve $\Fhat$ if it does so by
corrupting the conditional samples.

\paragraph{Conditional $Y\mid X$ diagnostics.}
\label{sec:conditional}
At a set of $x$-bin centres we compare the empirical conditional
$\phat(y\mid x\in I_j)$ (the $y$-histogram of particles whose $x$ falls in a
narrow window) to the reference conditional
$\pref(y\mid x)\propto e^{-\beta V(x,y)}$, via an $\Ltwo$ distance and a KL
divergence. By \cref{prop:cond_invariance} a benign FR step should leave these
close to the ABF-only values.

\subsection{GPU implementation and validation}
The production runs use a batched PyTorch backend that simulates many runs at
once and reuses the reference metrics/diagnostics code verbatim. For speed it
uses the binned-smoothed ABF/FR estimator of \cref{sec:abf_estimator} rather
than the exact kernel sums~\eqref{eq:nw}, and a fully vectorised fixed-$N$
birth--death. The two backends are not bit-identical (RNG and floating-point
order differ), but a validation suite passes all of its checks: the analytic
torch gradient matches finite differences and the numpy gradient; numpy-CPU
versus torch-CPU and torch-CPU versus torch-CUDA metrics agree to the same
order of magnitude; batch-size invariance holds; the binned-smooth estimator
matches the kernel reference to within the stated tolerance ($<10\%$ static,
well under $1\%$ in practice); and the FR event-fraction cap is respected. We
flag in \cref{sec:discussion} that this binned estimator is a deliberate, but
not mathematically identical, surrogate for the kernel notebook implementation.

\subsection{Hyperparameter grid and study size}
The FR methods sweep the rate $\gamma$, the marginal bandwidth $\eta$, and the
burn-in fraction $b$, at a fixed FR stride; ABF only has no FR parameters. The
resulting study (\cref{tab:design}) comprises \NRuns{} runs over \NConfigs{}
configurations and \NSeeds{} seeds, at \NParticles{} particles and \NSteps{}
steps each. We verified directly from the result CSVs that all \NRuns{} runs
completed with \emph{no} NaNs in the primary metrics. The selected
configurations reported in \cref{sec:results} are listed in
\cref{tab:selected_configs}.

\begin{table}[t]
\centering
\small
\caption{Production study design (read from the result CSVs and \texttt{configs/two\_dim\_xi\_x\_production\_gpu.yaml}).}
\label{tab:design}
\begin{tabular}{ll}
\toprule
Quantity & Value \\
\midrule
Inverse temperature $\beta$ & 4.0 \\
Time step $\Delta t$ & 0.002 \\
Steps per run & 100{,}000 \\
Particles per run $N$ & 1000 \\
Seeds & 5 (0, 1, 2, 3, 4) \\
Methods & 4 (ABF; FR estimated / uniform / oracle) \\
FR rate $\gamma$ & 0.01, 0.02, 0.05, 0.1 \\
FR bandwidth $\eta$ & 0.075, 0.1, 0.15 \\
FR burn-in fraction & 0, 0.2, 0.4 \\
FR stride \texttt{fr\_every} & 5 \\
Total configurations & 109 \\
Total runs & 545 \\
\bottomrule
\end{tabular}
\end{table}

\begin{table}[t]
\centering
\small
\caption{Hyperparameters of the selected configurations (integrated-$L^2(F)$ rule). $n_{\rm FR}$ is the Fisher--Rao stride \texttt{fr\_every}.}
\label{tab:selected_configs}
\begin{tabular}{llrrrrr}
\toprule
Method & target & $\gamma$ & $\eta$ & burn-in & $n_{\rm FR}$ & seeds \\
\midrule
ABF only & none & 0 & 0.075 & 0 & 1 & 5 \\
FR estimated & estimated & 0.1 & 0.075 & 0 & 5 & 5 \\
FR uniform & uniform & 0.05 & 0.15 & 0.2 & 5 & 5 \\
FR oracle* & oracle & 0.1 & 0.075 & 0 & 5 & 5 \\
\bottomrule
\end{tabular}
\end{table}

